\section{PROCEDIMENTOS METODOLÓGICOS}
\label{sec:metodos}

Trata-se de uma pesquisa de abordagem quantitativa e, quanto aos objetivos, descritiva, na acepção de \textcite{gil2017}: busca descrever as características de uma população (os municípios maranhenses) e estabelecer relações entre variáveis, sem manipulá-las. Quanto aos procedimentos técnicos, classifica-se como pesquisa documental \parencite{gil2017}, pois utiliza exclusivamente dados secundários de acesso público, produzidos pelo INEP e pelo Instituto Brasileiro de Geografia e Estatística (IBGE) com outras finalidades e aqui reanalisados. Esta seção apresenta as fontes, as variáveis e o tratamento dos dados e, em seguida, os dois estágios da análise.

\subsection{Fontes de dados e variáveis}

O estudo integra três bases públicas. A primeira é o IDEB 2025 dos anos iniciais do ensino fundamental (EFAI) da rede municipal \parencite{inep2025}, calculado conforme \textcite{fernandes2007}. A edição cobriu 216 dos 217 municípios do estado; Nova Iorque não teve o indicador divulgado porque o INEP não divulgou as notas do município no Saeb. A segunda são os microdados do Censo Escolar 2025 \parencite{inep2026microdados}, com a infraestrutura, as matrículas e as turmas de cada escola, complementados pela Tabela 2.29 da Sinopse Estatística da Educação Básica 2025 \parencite{inep2026sinopse}, que traz a formação dos docentes dos anos iniciais sem a dupla contagem presente nos microdados brutos. A terceira é o Censo Demográfico 2022, do qual foram extraídos, no Sistema IBGE de Recuperação Automática (Sidra), o rendimento médio domiciliar \textit{per capita} (tabela 10295) \parencite{ibge2022} e a distribuição dos moradores por classes de rendimento (tabela 10296) \parencite{ibge2022classes}. Por ser um recenseamento, é a única fonte de renda com cobertura de todos os municípios.

A segmentação usa duas variáveis do Censo 2022: a \textit{renda per capita}, que mede o nível de recursos do município, e a \textit{pobreza extrema}, percentual de moradores com rendimento domiciliar \textit{per capita} de até 1/4 do salário mínimo ou sem rendimento, critério de miserabilidade da Lei Orgânica da Assistência Social \parencite{brasil1993}. A escolha tem base teórica: o nível socioeconômico é o preditor mais consistente do desempenho desde \textcite{coleman1966}, e a associação é mais forte em unidades agregadas, como escolas e distritos, do que no nível individual \parencite{sirin2005}. Controlar o contexto é, portanto, condição para não atribuir a fatores escolares diferenças que decorrem da renda \parencite{albernaz2002}. As duas medidas se complementam, pois a média é sensível a rendas altas concentradas em poucos domicílios, e a pobreza extrema descreve a base da distribuição. Elas são correlacionadas (Pearson de $-0{,}87$), e a contribuição de cada uma é avaliada por ablação no Estágio 1.

As variáveis explicativas, organizadas em três grupos, são descritas no quadro~\ref{qua:variaveis}.

\begin{quadro}[!ht]
\centering
\caption{Variáveis explicativas utilizadas na análise intragrupo}
\label{qua:variaveis}
\corpodez
\begin{tabular}{|>{\raggedright\arraybackslash}p{0.42\textwidth}|>{\raggedright\arraybackslash}p{0.50\textwidth}|}
\hline
\multicolumn{1}{|c|}{\textbf{Variável}} & \multicolumn{1}{c|}{\textbf{Definição operacional}} \\
\hline
\multicolumn{2}{|>{\raggedright\arraybackslash}p{0.92\textwidth}|}{\textit{Infraestrutura escolar (Censo Escolar 2025): proporção de escolas com o atributo}} \\
\hline
\% escolas com água adequada & Água potável ou rede pública \\
\hline
\% escolas com biblioteca/sala de leitura & Biblioteca, biblioteca-sala de leitura ou sala de leitura \\
\hline
\% escolas com esgoto adequado & Rede pública ou fossa séptica \\
\hline
\% escolas com internet & Acesso à internet \\
\hline
\% escolas com laboratório de informática & Laboratório de informática \\
\hline
\% escolas com quadra esportiva & Quadra de esportes \\
\hline
\multicolumn{2}{|>{\raggedright\arraybackslash}p{0.92\textwidth}|}{\textit{Qualificação docente (Sinopse Estatística 2025)}} \\
\hline
\% docentes EFAI com ensino superior & Docentes dos anos iniciais com graduação \\
\hline
\% docentes EFAI com licenciatura & Docentes dos anos iniciais com licenciatura \\
\hline
\multicolumn{2}{|>{\raggedright\arraybackslash}p{0.92\textwidth}|}{\textit{Fluxo e distorção (Censo Escolar 2025)}} \\
\hline
Alunos por turma (educação básica) & Matrículas por turma na educação básica \\
\hline
Alunos por turma (ensino fundamental) & Matrículas por turma no ensino fundamental \\
\hline
Proxy distorção idade-série & \% de matrículas com 15 anos ou mais \\
\hline
\% turmas com correção de fluxo (anos finais) & Turmas de correção de fluxo entre as dos anos finais \\
\hline
\end{tabular}
\fonte{elaboração própria (2026).}
\end{quadro}

\subsection{Coleta e pré-processamento dos dados}

A coleta e o tratamento foram automatizados em \textit{scripts} Python, o que permite reproduzir a análise. As variáveis escolares da rede municipal foram agregadas por município como proporções de escolas, de docentes ou de turmas, ou como médias sobre as matrículas. A unidade de análise é, portanto, o município: as associações encontradas referem-se às redes municipais e não podem ser transpostas para escolas, professores ou alunos \parencite{robinson1950}.

A Sinopse conta os docentes dos anos iniciais de todas as redes, enquanto o IDEB analisado é o da rede municipal, que reúne, em média, 93\% (\textit{cluster} intermediário) e 95\% (vulnerável) desses docentes. Para verificar se o descompasso afeta os resultados, a formação docente também foi calculada apenas na rede municipal, a partir dos microdados, abrangendo toda a educação básica, porque os microdados não separam a formação por etapa. Integradas as bases, não há valores ausentes além do IDEB de Nova Iorque, que, sem poder ser classificado por desempenho, fica fora dos dois estágios.

\subsection{Estágio 1: segmentação socioeconômica por \textit{K-Means}}

O \textit{K-Means} \parencite{macqueen1967} agrupa os 216 municípios pela renda \textit{per capita} e pela pobreza extrema, padronizadas pela mediana e pelo intervalo interquartil para reduzir a influência de municípios atípicos, como São Luís e Imperatriz, cujos valores extremos de renda e de pobreza poderiam distorcer a formação dos grupos.

O algoritmo é executado para um número de grupos $k$ entre 2 e 10, e $k$ é escolhido pelo \textit{Silhouette Score} médio \parencite{rousseeuw1987}. A escolha é confirmada pela estabilidade da partição, avaliada por \textit{bootstrap} de 25 iterações com 80\% dos dados, na configuração de referência de \textcite{benhur2002}, com o \textit{Adjusted Rand Index} (ARI) como métrica de concordância. Na ablação, a clusterização é refeita sem cada variável, e o ARI entre a nova partição e a original indica o quanto a partição depende dela. Os \textit{clusters} são visualizados por Análise de Componentes Principais (\textit{Principal Component Analysis}, PCA) \parencite{abdi2010} com duas componentes.

\subsection{Estágio 2: análise intragrupo}

Em cada \textit{cluster} com amostra suficiente, comparam-se os municípios de \textit{alto IDEB} (quartil superior, $\geq Q_{75}$) e de \textit{baixo IDEB} (quartil inferior, $\leq Q_{25}$); o intervalo central é excluído, o que aumenta o contraste e reduz o ruído.

Para cada variável, o teste de Mann-Whitney \parencite{mannwhitney1947} bicaudal verifica se as distribuições diferem entre os grupos, e o $\delta$ de Cliff \parencite{cliff1993}, de $-1$ a $1$, mede a magnitude da diferença. Seguindo \textcite{vargha2000}, valores absolutos a partir de 0,12, 0,28 e 0,42 são lidos como efeitos pequeno, médio e grande. Como são testadas catorze variáveis em cada \textit{cluster} (as doze do quadro~\ref{qua:variaveis} e as duas de segmentação), os p-valores são acompanhados do valor $q$ de \textcite{benjamini1995}, que controla a taxa de falsas descobertas. São reportadas as variáveis com $p < 0{,}10$, interpretadas à luz do valor $q$, do efeito e da consistência entre \textit{clusters}, o que, com o tamanho das amostras, dá ao estágio 2 caráter exploratório.

Duas verificações complementam a comparação: a correlação de Spearman entre cada variável e o IDEB com todos os municípios do \textit{cluster}, sem descartar o intervalo central; e o poder preditivo conjunto das variáveis, em validação cruzada em cinco partes repetida vinte vezes, de um classificador \textit{Random Forest} \parencite{breiman2001}, medido pela área sob a curva ROC (\textit{Receiver Operating Characteristic}), a AUC, em que 0,5 equivale ao acaso, e de uma regressão \textit{ridge} \parencite{hoerl1970} para o IDEB contínuo, medida pelo $R^2$ fora da amostra.

\subsection{Uso de ferramentas de inteligência artificial generativa}
\label{sec:uso_ia}

Foram utilizadas ferramentas de inteligência artificial generativa em três etapas: o Claude (Anthropic), na correção e revisão dos \textit{scripts} Python da análise, com cada alteração validada pela autoria; o Claude e o ChatGPT (OpenAI), na revisão da redação e da clareza do texto; e o Claude e o SciSpace AI Research (PubGenius Inc.), na busca de trabalhos relevantes e na verificação de metadados de citações, com cada referência conferida na fonte original. O texto, a argumentação e a interpretação dos resultados são de responsabilidade exclusiva da autoria, conforme as diretrizes de integridade científica do Conselho Nacional de Desenvolvimento Científico e Tecnológico (CNPq).
